\subsection{Optimal affine error for every cyclic type}

\begin{theorem}[DP43: the optimal affine loss is the codimension]\label{DP43}
For the arbitrary-unit cyclic display of DP40, let $m=|A|$, $\ell=h-m$, and define at every cyclic position $a$, selected or not,
\[
 c_a=\frac{\sum_{b\notin A}p^{[a-b]_h}}{q-1},\qquad
 r_a=\frac1{p-1}-c_a.
\]
For $i\in A$ put
\[
 E_i(n)=\sum_{\substack{u\ge1\\i+u\notin A}}
                  \floor{n/p^u},
\]
where membership is read modulo $h$. Then, coefficientwise for every integral $f$,
\begin{equation}\label{eq:cyclicaffine}
 \vp(D_i^nf)\ge
 \max\{E_i(n),\ceil{c_in-(h-m)}\}.
\end{equation}
The slope $c_i$ and the constant $h-m$ are both optimal for every fixed normalized eigendirection. If $g_a$ is the forward gap from $a\in A$ to its next selected index, the single monomial
\[
 f_A=\prod_{a\in A}X_a^{p^{g_a}-1}
\]
attains the affine bound in every direction:
\[
 \vp(D_i^{q-1}f_A(0))=c_i(q-1)-(h-m).
\]
All primes, including $p=2$, are covered. The floor sum and affine bound are two lower bounds; no exact optimum at every order is asserted.
\end{theorem}

\begin{proof}
Write $s=1/(p-1)$ and $\epsilon_a=1$ off $A$, zero on $A$. The weights at all cyclic positions, including the exponent-zero contribution off $A$, satisfy
\[
 0<c_a,r_a<s,\qquad
 c_a=\sum_{u\ge1}\epsilon_{a+u}p^{-u}.
\]
Consequently
\[
 p^tc_i=b_i(t)+c_{i+t},\qquad
 b_i(t)=\sum_{u=1}^t\epsilon_{i+u}p^{t-u}\in\Z_{\ge0}.
\]
DP40's contraction at $n=p^t$ gives the coefficientwise bound
\[
 \vp(D_i^{p^t}f)\ge p^tc_i-s=b_i(t)+(c_{i+t}-s).
\]
All finite coefficient valuations are integers because $R=W(k)$ is unramified, and $-1<c_{i+t}-s<0$, also at $p=2$.
Thus $D_i^{p^t}$ preserves $p^{b_i(t)}R[[\mathbf X]]$ in the sense that it sends every integral function into that ideal. Compose these operator inclusions for the base-$p$ digits $n=\sum_tn_tp^t$. Interchanging finite sums gives
\[
 \vp(D_i^nf)\ge\sum_tn_tb_i(t)
       =\sum_{u\ge1}\epsilon_{i+u}\floor{n/p^u}=E_i(n).
\]

\emph{A primitive reduction of the normalized $q$th derivative.}
Fix $i\in A$ and put
\[
 C_i=c_i(q-1)=
       \sum_{\substack{1\le u<h\\i+u\notin A}}p^{h-u}\in\Z_{>0}.
\]
Because $i+h\in A$, $b_i(h)=E_i(q)=C_i$. Therefore
$\mathscr T_i=p^{-C_i}D_i^q$ is an integral continuous operator commuting with translations.
For $0<k<q$, put $t=\vp(k)$.
The equality
\[
 \vp\binom qk=h-t
\]
follows from $\binom qk=(q/k)\binom{q-1}{k-1}$, since every coefficient of
$(1+Z)^{q-1}$ is a unit modulo $p$.
For $1\le u<h$ one has
\[
 \floor{k/p^u}+\floor{(q-k)/p^u}
       =q/p^u-\mathbf1_{u>t}.
\]
It follows that
\[
 E_i(k)+E_i(q-k)
       =C_i-\#\{u:t<u<h,\ i+u\notin A\}.
\]
Every cross term in the Leibniz expansion of $\mathscr T_i(fg)$ is therefore divisible by $p$, since its valuation is at least
\[
 h-t-\#\{u:t<u<h,\ i+u\notin A\}\ge1.
\]
Hence $\overline{\mathscr T}_i$ is an invariant derivation of the special-fibre formal group. An invariant derivation is determined by its tangent value at the identity, by differentiating the translation identity. It is therefore a residue-field linear combination of the $\bar D_a$ for $a\in A$.

\emph{The torus isolates the same direction.}
Make a finite unramified coefficient extension containing $\F_q$ and a primitive $(q-1)$st root $\zeta$. The display automorphism
\[
 e_a\longmapsto[\zeta^{p^a}]e_a
\]
commutes with its arrows, even for arbitrary $u_a$, because Witt Frobenius takes $[\zeta^{p^a}]$ to $[\zeta^{p^{a+1}}]$. It induces exactly $X_a\mapsto\zeta^{p^a}X_a$ on the Teichm\"uller coordinates.
The character of $D_a$ in
$D_aa_\zeta^*=\zeta^{p^a}a_\zeta^*D_a$ is $\zeta^{p^a}$.
The character of $\mathscr T_i$ is $\zeta^{qp^i}=\zeta^{p^i}$.
These $h$ characters are distinct also in the residue field. Consequently
$\overline{\mathscr T}_i=\bar a_i\bar D_i$.
DP40's exact $q$th coordinate coefficient makes
$a_i:=\mathscr T_iX_i(0)$ a unit already in the original ring.
The equality descends from the unramified extension, so
\begin{equation}\label{eq:cyclicqoperator}
 D_i^q=p^{C_i}a_iD_i+p^{C_i+1}\mathscr R_i,
\end{equation}
where $\mathscr R_i$ is an integral continuous $R$-linear operator. The entire operator identity, including the extra $p$, has now been proved; one sharp coordinate value alone would not suffice.

\emph{Affine induction.}
For $0\le n<q$ with $n=\sum_{t=0}^{h-1}n_tp^t$,
\[
 c_in-E_i(n)=\sum_{t=0}^{h-1}n_tc_{i+t}
       \le(p-1)\sum_{a\in\Z/h\Z}c_a=\ell.
\]
The last equality counts each of the $\ell$ complementary indices against
$\sum_{a=0}^{h-1}p^a=(q-1)/(p-1)$.
For $n\ge q$, multiply \eqref{eq:cyclicqoperator} by $D_i^{n-q}$ and induct. The valuations of its two terms are at least
\[
 C_i+c_i(n-q+1)-\ell=c_in-\ell,\qquad
 C_i+1+c_i(n-q)-\ell=c_in-\ell+1-c_i.
\]
Since $c_i<s\le1$ and $\mathscr R_i$ preserves integral coefficient ideals, the induction holds. Integer rounding and the independently proved floor sum give \eqref{eq:cyclicaffine}.

\emph{A common sharp witness.}
Let $e_a=p^{g_a}-1$. The selected-edge weight identities give
$p^{g_a}r_a-r_{\operatorname{next}(a)}=1$; summing yields
\[
 \sum_{a\in A}e_ar_a=m.
\]
Fix an input index $i\in A$. If $d(i,a)$ is the forward distance from $i$ to $a$, set $N_a=p^{d(i,a)}$.
The first nonzero pure-$t$ term of $E_a(te_i)$ has degree exactly $N_a$ and coefficient valuation $r_a-N_ar_i$.
Indeed torus equivariance forces $Np^i\equiv p^a\bmod(q-1)$, whose least positive solution is $N_a$.
In DP40's contracted inverse $\sum_{b\ge0}(-1)^bM_\eta^{\circ b}$, the unique path before a complete cycle supplies a nonzero coefficient at that degree.
These two observations prove both minimality and exact valuation, with no cancellation from unit arrows.

List the selected indices cyclically from $i$. Multiplication of $N_a$ by $p^{g_a}$ gives the next degree, with the last giving $q$; thus
$\sum_{a\in A}e_aN_a=q-1$.
The first coefficient of $f_A(E(te_i))$ occurs at that degree and is the product of the corresponding first coefficients, with valuation
\[
 \sum_{a\in A}e_a(r_a-N_ar_i)=m-(q-1)r_i.
\]
Finally $\vp((q-1)!)=s(q-1)-h$, so the factorial derivative has valuation
$c_i(q-1)-(h-m)$, as asserted.
A smaller real affine constant would contradict the ceiling bound at this order, because $c_i(q-1)$ is an integer. The slope itself is sharp by DP40's $q^r$ sequence.
\end{proof}

The singleton type gives the affine loss $h-1$ of DP1, while its complementary type gives loss one as in DP34. Mixed derivative estimates follow by composing the coefficientwise operator inclusions. This codimension statement belongs to the integral display and distribution lattice; it does not follow merely from the rational Newton slopes of an arbitrary deformation.

